\section{Conclusion}

\label{sec:conclusion}

Nebula extends the metastable consensus family to sparse and unreliable networks
through stratified sampling with bridge validators. The expansion factor framework
provides a clean characterization of how network quality affects consensus performance,
with graceful degradation from $O(\log n)$ in well-connected networks to $O(n \log n)$
at the connectivity threshold. Nebula enables consensus in environments where standard
protocols fail: high packet loss, validator churn, temporary partitions, and small
appchain validator sets.

\begin{thebibliography}{99}
\bibitem{rocket2019wave} T. Rocket et al., ``Scalable and probabilistic leaderless BFT consensus through metastability,'' 2019.
\bibitem{kelling2020wave} Z. Kelling, ``Wave protocol: Decision stability and confidence monotonicity,'' Lux Partners Limited Limited, 2020.
\bibitem{kelling2023nova} Z. Kelling, ``Nova protocol: Supernova burst consensus for high-throughput chains,'' Lux Partners Limited Limited, 2023.
\bibitem{penrose2003random} M. Penrose, \emph{Random Geometric Graphs}, Oxford University Press, 2003.
\bibitem{dwork1988consensus} C. Dwork, N. Lynch, and L. Stockmeyer, ``Consensus in the presence of partial synchrony,'' \emph{JACM}, 1988.
\bibitem{gilbert2002brewer} S. Gilbert and N. Lynch, ``Brewer's conjecture and the feasibility of consistent, available, partition-tolerant web services,'' \emph{SIGACT News}, 2002.
\end{thebibliography}

